% ==============================================================================
% secoes/00_resumo.tex - RESUMO E PALAVRAS-CHAVE
% Tema: Federico Faggin e o Primeiro Microprocessador Monolítico (Intel 4004)
% ==============================================================================

\begin{abstract}
% [DIRETIVA DE ESCRITA: RESUMO (ABSTRACT)]
% - Extensão obrigatória: Máximo de 250 palavras (conforme Orientacoes_Trabalho.pdf).
% - Caráter: Autoexplicativo, conciso e redigido em voz ativa e tempo pretérito/presente.
% - Conteúdo mandatório a cobrir sinteticamente:
%   1. Problema de Engenharia: A barreira de interconexão e a inviabilidade de integrar
%      uma CPU em chip único com tecnologia convencional de porta metálica de alumínio.
%   2. Metodologia: Abordagem qualitativa de pesquisa histórico-documental,
%      triangulando patentes de Faggin (Silicon Gate Technology - SGT), notas de projeto
%      e manuais originais do Intel 4004 (MCS-4) com literatura de educação em engenharia.
%   3. Resultados Principais: Como a tecnologia SGT (1968), as células dinâmicas com bootstrapping
%      e o roteamento bidimensional silício-metal permitiram integrar 2.300 transistores
%      em 12 mm², viabilizando a CPU monolítica e o paradigma de firmware/sistemas embarcados.
%   4. Discussão e Conclusão Pedagógica: Articulação entre física de semicondutores,
%      arquitetura de computadores, linguagens HDL e co-design hardware/software no ensino contemporâneo.

Este artigo analisa a contribuição pioneira do físico e engenheiro ítalo-americano Federico Faggin (1941--) para o surgimento do primeiro microprocessador comercial em uma única pastilha de silício, o Intel 4004 (1971), e avalia seus impactos na sociedade moderna e na formação em engenharia. Partindo de uma abordagem histórico-pedagógica, examina-se como o desenvolvimento da tecnologia MOS com porta autoalinhada de silício (\textit{Silicon Gate Technology} -- SGT) e as técnicas avançadas de projeto lógico e layout viabilizaram a integração de uma unidade central de processamento (CPU) completa em um chip monolítico. A metodologia fundamenta-se em pesquisa documental de fontes primárias (patentes, notas técnicas originais e manuais do Intel 4004) e secundárias, trianguladas com literatura em educação em engenharia. Os resultados demonstram que o 4004 não apenas inaugurou a era dos sistemas computacionais embarcados e da automação digital ubíqua, mas também unificou os domínios de dispositivos semicondutores e arquitetura de computadores. Discute-se a relevância pedagógica de utilizar a gênese do microprocessador para articular a física do estado sólido, linguagens de descrição de hardware (HDLs) e o paradigma de co-design hardware/software nos currículos contemporâneos. Conclui-se que o legado de Faggin exemplifica a síntese entre engenharia de processos e visão sistêmica, oferecendo um modelo formativo indispensável para futuros engenheiros.
\end{abstract}

\begin{IEEEkeywords}
% [DIRETIVA DE ESCRITA: PALAVRAS-CHAVE (KEYWORDS)]
% - Quantidade obrigatória: Entre 4 e 6 termos (conforme Orientacoes_Trabalho.pdf).
Federico Faggin, Intel 4004, Silicon Gate Technology, Microprocessadores, Sistemas VLSI, Educação em Engenharia.
\end{IEEEkeywords}
